\documentclass[11pt,a4paper]{article}
\usepackage[utf8]{inputenc}
\usepackage[T1]{fontenc}
\usepackage[spanish,es-tabla]{babel}
\usepackage{geometry}
\geometry{top=2.5cm, bottom=2.5cm, left=2.5cm, right=2.5cm}
\usepackage{amsmath,amssymb}
\usepackage{listings}
\usepackage{xcolor}
\usepackage{hyperref}
\usepackage{booktabs}
\usepackage{fancyhdr}

\lstset{
  language=Caml,
  basicstyle=\small\ttfamily,
  keywordstyle=\color{blue}\bfseries,
  commentstyle=\color{gray}\itshape,
  stringstyle=\color{red},
  numbers=left,
  numberstyle=\tiny\color{gray},
  stepnumber=1,
  backgroundcolor=\color{black!3},
  frame=single,
  rulecolor=\color{black!15},
  tabsize=2,
  breaklines=true,
  breakatwhitespace=true
}

\pagestyle{fancy}
\fancyhf{}
\fancyhead[L]{Paradigmas de Programación}
\fancyhead[R]{Recopilación de Enunciados}
\fancyfoot[C]{\thepage}

\title{\textbf{Recopilación de Enunciados de Prácticas}\\Paradigmas de Programación (OCaml)}
\author{Grado en Ingeniería Informática}
\date{\today}

\begin{document}

\maketitle
\tableofcontents
\newpage

\section{Práctica 4: Conjetura de Collatz y Algoritmo de Euclides}
\subsection{Ejercicio 1: Conjetura de Collatz}
Dada la función $f(n)$ definida en OCaml como:
\begin{lstlisting}
let f n = if n mod 2 = 0 then n / 2 else 3 * n + 1
\end{lstlisting}
La Conjetura de Collatz afirma que, comenzando con cualquier número entero positivo, la secuencia siempre alcanzará el número 1. Se solicita implementar en el archivo \texttt{collatz.ml} las siguientes funciones:
\begin{itemize}
    \item \texttt{check\_to : int -> bool}: Comprueba la conjetura para todos los números naturales menores o iguales que el argumento dado.
    \item \texttt{orbit : int -> string}: Devuelve una cadena con la secuencia de elementos de la órbita del número, separados por comas y un espacio.
    \item \texttt{length : int -> int}: Devuelve la longitud de la órbita de un número (número de pasos para llegar a 1).
    \item \texttt{top : int -> int}: Devuelve el valor máximo alcanzado en la órbita del número.
\end{itemize}

\subsection{Ejercicio 2 (Opcional): Optimización de Collatz}
Implemente en \texttt{collatz\_plus.ml} una versión optimizada que calcule de forma simultánea el par de la longitud y el valor máximo de la órbita en un único recorrido:
\begin{lstlisting}
val length_and_top : int -> int * int
\end{lstlisting}

\subsection{Ejercicio 3: Algoritmo de Euclides}
Implemente en \texttt{mcd.ml} dos versiones de la función para calcular el Máximo Común Divisor (MCD) de un par de números enteros no negativos:
\begin{itemize}
    \item \texttt{mcd : int * int -> int}: Versión clásica restando sucesivamente de manera recursiva.
    \item \texttt{mcd' : int * int -> int}: Versión optimizada usando el resto de la división entera (\texttt{mod}).
\end{itemize}


\section{Práctica 5: Sucesión de Fibonacci y Operaciones de Consola}
\subsection{Ejercicio 1: Programa \texttt{fibto}}
Escriba un programa ejecutable en un archivo \texttt{fibto.ml} de modo que su ejecución con un argumento $n$ muestre por la salida estándar todos los términos de la sucesión de Fibonacci menores que $n$.
\begin{itemize}
    \item Debe validar los argumentos usando \texttt{Sys.argv}.
    \item Se prohíbe explícitamente el uso de bucles imperativos como \texttt{while} y \texttt{for}, debiendo utilizarse recursión.
    \item No está permitido el uso del punto y coma (\texttt{;}) secuencial ni el doble punto y coma (\texttt{;;}).
\end{itemize}

\subsection{Ejercicio 2: Fibonacci Eficiente}
Optimice la función de cálculo en \texttt{fastfibto.ml} para permitir el procesamiento inmediato de números grandes evitando la redundancia exponencial de la implementación ingenua.

\subsection{Ejercicio 3, 4 y 5 (Opcionales): Conversor de Bases}
En \texttt{strg.ml}, implemente funciones de conversión de enteros a cadenas en diferentes bases numéricas:
\begin{itemize}
    \item \texttt{str3\_of\_int : int -> string} y su inversa \texttt{int\_of\_str3 : string -> int} (base 3).
    \item \texttt{strg\_of\_int : int -> int -> string} y su inversa \texttt{int\_of\_strg : int -> string -> int} (bases genéricas de 2 a 9).
\end{itemize}


\section{Práctica 6: Torres de Hanoi y Teoría de Números}
\subsection{Ejercicio 1 y 2: Torres de Hanoi}
Dada la implementación incompleta de las Torres de Hanoi, complete \texttt{hanoi.ml} reconstruyendo las funciones:
\begin{lstlisting}
val mueve : int * int -> string
val otro : int -> int -> int
\end{lstlisting}
Y añada la función \texttt{n\_hanoi\_mov : int -> int -> int -> int -> int * int} que devuelve directamente el $n$-ésimo movimiento de la secuencia óptima para resolver el problema con $nd$ discos, con una complejidad temporal de $\mathcal{O}(nd)$.

\subsection{Ejercicio 3 y 4: Funciones sobre Primos y Conjetura de Goldbach}
En \texttt{prime.ml}, implemente:
\begin{itemize}
    \item \texttt{is\_prime2 : int -> bool}: Versión optimizada para determinar si un número es primo.
    \item \texttt{goldbach : int -> int * int}: Devuelve un par de números primos que sumen el entero par dado, según la Conjetura de Goldbach.
\end{itemize}

\subsection{Ejercicio 5 (Opcional): Currificación y Composición}
Implemente en \texttt{ej65.ml} funciones para la currificación (\texttt{curry}), descurrificación (\texttt{uncurry}) y composición de funciones (\texttt{comp}).


\section{Práctica 7: Excepciones y Módulo de Listas}
\subsection{Ejercicio 1: Factorial Robusto}
Redefina la función factorial en \texttt{fact.ml} para que, al aplicarse sobre números negativos, lance la excepción \texttt{Invalid\_argument "fact"} asegurándose de realizar la comprobación de signo únicamente una vez.

\subsection{Ejercicio 2: Reimplementación de myList}
Escriba \texttt{myList.ml} con las funciones clásicas del módulo \texttt{List} sin usar el operador de concatenación original (\texttt{@}). Implemente versiones directas e interfaces recursivas terminales (ej. \texttt{append'}, \texttt{map'}, \texttt{filter'}).


\section{Práctica 8: Iteración con \texttt{fold\_left}}
\subsection{Ejercicio 1: Módulo folding}
Implemente \texttt{folding.ml} usando \textbf{únicamente} \texttt{List.fold\_left}. Se prohíbe el uso de recursión explícita o cualquier otra función del módulo \texttt{List}.

\subsection{Ejercicio 2: Módulo listing}
Implemente \texttt{listing.ml} usando funciones del módulo \texttt{List} de manera que toda la recursión sea terminal.


\section{Práctica 9: Concatenación Eficiente y Tail Modulo Constructor}
\subsection{Ejercicio 1: Concatenación Terminal}
Explique por qué acumular listas con \texttt{List.append} dentro de un \texttt{List.fold\_left} genera problemas de rendimiento ($\mathcal{O}(N^2)$). Escriba en \texttt{concat.ml} las variantes eficientes \texttt{concat'} (con \texttt{fold\_right}) y \texttt{concat} (recursiva terminal).

\subsection{Ejercicio 2: Sublistas}
Añada \texttt{sublists : 'a list -> 'a list list} para generar el conjunto potencia de una lista manteniendo el orden relativo original de los elementos.

\subsection{Ejercicio 3 y 4: Tail Modulo Constructor (TMC)}
Defina versiones recursivas terminales tradicionales en \texttt{tail.ml} y explore la optimización de compilador \textit{Tail Modulo Constructor} en \texttt{tail2.ml}.


\section{Práctica 10: Movimientos en Cuadrícula (El Patio de la Cárcel)}
\subsection{Ejercicio 1: Escape de Prisión (\texttt{tour.ml})}
Dadas las dimensiones de un patio de $m \times n$ baldosas con robots patrullando que rebotan en los bordes, escriba la función:
\begin{lstlisting}
val tour : int -> int -> ((int * int) * move) list -> move list
\end{lstlisting}
que determine un camino libre de colisiones simultáneas para que el preso escape desde $(1,1)$ hasta $(m,n)$.

\subsection{Ejercicio 2 (Opcional): Ruta Óptima (\texttt{shortest.ml})}
Implemente la función \texttt{shortest} para calcular la ruta de escape de longitud mínima.


\section{Práctica 11: Árboles Binarios, Estrictos, Generales y BST}
\subsection{Ejercicio 1: Recorrido de BinTree}
Construya el módulo \texttt{BinTree} con una función de recorrido en anchura (\texttt{breadth}) optimizada para evitar la ineficiencia de la concatenación tradicional de listas.

\subsection{Ejercicio 2 y 3: Árboles Estrictamente Binarios y Árboles Multicamino}
Implemente \texttt{StBinTree.ml} y \texttt{gTree.ml} para trabajar con variantes estructurales de árboles.

\subsection{Ejercicio 4 (Opcional): Árboles Binarios de Búsqueda}
Implemente \texttt{bst.ml} con funciones \texttt{is\_bst}, \texttt{mem} (iterativa sin recursividad mediante bucles \texttt{while} y referencias), \texttt{add} y \texttt{remove}.


\section{Práctica 12: Intérprete Simple C y Árboles en Arrays}
\subsection{Ejercicio 1 (Opcional): Intérprete de Simple C}
Escriba en \texttt{simpleC.ml} las correcciones del intérprete para las estructuras de control \texttt{While} y \texttt{For}. Incorpore controles de integridad y la función recursiva doble \texttt{legal\_body}.

\subsection{Ejercicio 2: Árboles en Array}
En \texttt{aTree.ml}, defina funciones para serializar y consultar árboles binarios en un array utilizando un tipo \texttt{'a option array}. Implemente \texttt{from\_bin}, \texttt{breadth} y \texttt{mem} (iterativa sin recursividad).

\subsection{Ejercicio 3 (Opcional): Colas y Bicolas}
En \texttt{bq.ml}, modifique el recorrido en anchura de los árboles utilizando un objeto de tipo cola eficiente, y extienda el comportamiento para definir una bicola (\texttt{queue'}).

\end{document}
